\section{Additional Related Work}
\label{sec:additional-related-work}

\subsection{Benchmarks Not Suitable for Measuring Self-Evolution}
\label{sec:additional-related-work-not-self-evolution}

Many benchmarks evaluate whether models or agents can answer questions, complete tasks, or produce correct artifacts, without making self-evolution itself the object of measurement. We use this label in a narrow evaluation sense: these benchmarks may still involve reasoning, tools, iteration, or environment interaction, but their primary reported quantity is not related to learning. It is usually final-answer accuracy, task success, pass rate, artifact quality, reproduction fidelity, or human-time horizon. In such settings, interaction is typically a means to produce a final answer or artifact, rather than the quantity being measured.

At the shortest and least interactive end, many classic capability benchmarks, such as MMLU~\cite{hendrycks2021mmlu}, GPQA~\cite{rein2023gpqa}, AIME~\cite{maa2026aime}, and closed-form coding or math evaluations~\cite{chen2021humaneval,hendrycks2021math}, ask models to solve static problems and report final accuracy. These benchmarks can be difficult and useful, but the model's actions do not create a changing environment, and feedback from the benchmark is not exposed as experience for later adaptation. They therefore measure static knowledge and reasoning accuracy rather than self-evolution.

Agentic software benchmarks increase realism by placing agents in codebases, development workflows, or larger construction tasks, but many still reduce evaluation to the final result. SWE-bench~\cite{jimenez2024swebench} evaluates issue repair in existing repositories; development and evolution benchmarks such as RoadmapBench~\cite{xu2026roadmapbench} and SWE-EVO~\cite{le2025sweevo} evaluate incremental changes over existing projects or release histories; construction or reconstruction benchmarks such as NL2Repo-Bench~\cite{ding2025nl2repobench} and ProgramBench~\cite{yang2026programbench} evaluate complete repository generation or recovery of program behavior; and SWE-AGI~\cite{zhang2026sweagi} evaluates specification-driven system construction under a fixed MoonBit API scaffold. FrontierCode~\cite{cognition2026frontiercode} raises the bar on readiness for production by evaluating whether coding agents produce maintainable, mergeable changes under repository-specific quality rubrics. Their task formats differ, but the shared evaluation target is final patch correctness, repository quality, program behavior, maintainability, or test pass rate rather than the trajectory by which an agent improves from feedback.

Some evaluations broaden realism without becoming agentic learning benchmarks. GDPval~\cite{patwardhan2025gdpval}, for example, measures model performance on economically valuable professional deliverables across occupations. Agents' Last Exam~\cite{sun2026agentslastexam} is a close comparator in realism and autonomy: it evaluates generalist computer-use agents on economically valuable professional workflows in real Windows or Linux sandboxes, with CLI and GUI access, professional software, hidden references, and verifiable final artifacts. The distinction is the evaluation target rather than agenticity. Agents' Last Exam is oriented toward completion: agents work toward a final deliverable, and the benchmark reports final scores, pass rates, cost, or time after the artifact is graded. Its trajectories are valuable for replay, audit, and failure analysis, but learning from the benchmark's feedback over a run is not the main quantity being measured.

Other benchmarks use realistic or long-horizon workflows that could support learning, but their published protocols mainly measure final outcomes. HCAST~\cite{rein2025hcast} and METR time-horizon evaluations~\cite{kwa2025timehorizon} calibrate model or agent success by human task duration. Terminal-Bench~\cite{merrill2026terminalbench}, RE-Bench~\cite{wijk2024rebench}, PaperBench~\cite{starace2025paperbench}, CORE-Bench~\cite{siegel2024corebench}, PRBench~\cite{qiu2026prbench}, ReplicationBench~\cite{ye2025replicationbench}, Collider-Bench~\cite{faroughy2026colliderbench}, and ScienceAgentBench~\cite{chen2024scienceagentbench} involve executable software tasks, terminal workflows, or scientific reproduction. These settings are more compatible with learning than short closed-form tasks, but their central questions are typically whether a task is completed, an artifact works, a paper is reproduced, or a final score is high enough under a budget. PaperBench is a useful example of the distinction: it evaluates agentic replication of 20 ICML papers from scratch using detailed rubrics and many gradable subtasks, but the reported score is still replication quality rather than improvement per unit of feedback.

\subsection{Benchmarks Suitable for Measuring Learning or Self-Evolution}
\label{sec:additional-related-work-learning}

Another line of work is more suitable for measuring learning or self-evolution because it exposes new context, serialized streams, iterative feedback, or long executable workspaces. These benchmarks are the closest conceptual comparisons to \benchmark{}, but they differ in what creates the experience stream and what quantity is ultimately scored. We group them by evaluation interface: learning from supplied context, learning over serialized streams or task sequences, and iterative optimization. This grouping is descriptive rather than an ordering by interaction strength. Sequential benchmarks can be highly interactive, and optimization benchmarks can expose diagnostic feedback; the key question is whether the agent's own behavior shapes the future experience it receives within a long executable task.

Context-learning benchmarks study the least agentic form of adaptation. CL-bench~\cite{dou2026clbench} and CL-bench Life~\cite{dou2026clbenchlife} test whether models can use newly provided professional or personal context to answer downstream questions or perform tasks grounded in that context. This counts as learning from external information, but the stream of episodes is usually not shaped by the agent's own actions, and there is no changing environment. The model is asked to use context, not to discover and exploit feedback through sustained interaction.

Sequential learning benchmarks introduce a stream-like structure. EvaLearn~\cite{dou2025evalearn} groups 648 problems into 182 sequences and evaluates learning capability and efficiency as models solve related problems in order. Continual Learning Bench~\cite{asawa2026clbench} explicitly reports improvement over sequential experience and introduces a stateful-versus-stateless comparison similar to Section~\ref{sec:stateful-stateless-analysis}. However, its synthetic task sequences have explicit subtask boundaries, whereas \benchmark{} uses single continuous problems. We therefore partition by time rather than by subtask, which changes how the stateful and stateless settings are defined. StreamBench~\cite{wu2024streambench} also reports improvement over feedback streams. LLF-Bench~\cite{cheng2023llfbench}, LifelongAgentBench~\cite{zheng2025lifelongagentbench}, and Evo-Memory~\cite{wei2025evomemory} study related questions through language feedback, interdependent tasks, experience replay, or evolving memory. These works establish that static capability and learning ability are distinct axes. The difference from \benchmark{} is not that their interaction is necessarily weaker; rather, the sequence is usually organized by the benchmark as a series of instances, tasks, feedback events, or memory updates. In \benchmark{}, the sequence is endogenous to one long task: the agent plans, edits artifacts, submits attempts, receives diagnostics, and thereby affects what it learns next.

Iterative optimization benchmarks are closer to \benchmark{} because they make repeated attempts and empirical feedback central to performance. Frontier-Eng~\cite{chi2026frontiereng} studies improvement over generative optimization cycles, while MLS-Bench~\cite{lyu2026mlsbench} and Frontier-CS~\cite{mang2025frontiercs} evaluate test-time discovery or open-ended CS optimization with visible metrics, submissions, or trajectory-level reporting. MLE-bench~\cite{chan2024mlebench} is also naturally viewed in this family: agents work in a free Kaggle-style workspace, the paper includes studies of how performance scales with resources, and its 100-hour analysis grades snapshots of the agent's best attempt over elapsed time. ALE-Bench~\cite{imajuku2025alebench} is another adjacent long-horizon setting driven by fixed objectives, even though its original protocol primarily reports outcomes on algorithm engineering rather than learning metrics.

FrontierSWE~\cite{chu2026frontierswe} and AutoLab~\cite{xu2026autolab} are the closest prior work. Both evaluate long-horizon agentic improvement over executable artifacts through repeated edits, experiments, and empirical feedback. FrontierSWE focuses on software engineering and performance-tuning tasks, with reported average agent runtime of roughly 3--4 hours across models and task categories. AutoLab gives agents working but deliberately suboptimal programs across systems, CUDA, model development, and puzzle-style optimization tasks; it directly measures iterative improvement and should not be treated as a non-learning benchmark, although performance curves are not reported as the main result and most current tasks run for 2--4 hours, with only a small minority exceeding 6 hours. Under the shared premise of long-horizon agentic tasks, the main distinction is domain coverage: FrontierSWE is concentrated on software engineering, AutoLab emphasizes research and engineering tasks centered on optimization, and \benchmark{} spans a broader set of executable domains. \benchmark{} also uses a day-scale task contract and makes the time-aligned trajectory itself the primary object of measurement, with metrics for improvement area, regression, and active learning span.

Taken together, these benchmarks cover important pieces of the problem: realistic executable work, learning over organized streams, and iterative optimization under feedback. \benchmark{} targets their intersection. It evaluates within-run self-evolution in long-horizon executable environments where the experience stream arises from the task itself, the agent can influence what it observes next, and evaluation uses a general-purpose agent harness rather than a benchmark-specific learning scaffold. Its trajectory-level metrics report improvement, regression, active learning span, and final performance over the same continuous run.

\subsection{Scaling Laws for LLMs and Agents}
\label{sec:additional-related-work-scaling-laws}

Classical neural and language-model scaling laws mainly study pretraining, modeling reducible loss as a function of model size, data, and training compute. Kaplan et al.~\cite{kaplan2020scaling} showed that language-modeling loss follows power-law relationships over several axes of scale, and Chinchilla-style work refined the compute-optimal allocation between model parameters and training tokens~\cite{hoffmann2022training}. More recent work studies how benchmark performance, rather than loss, changes with scale; because accuracy-like metrics are bounded, these curves are often better described by sigmoidal or other saturating forms~\cite{owen2024predictable,ruan2024observational,bhagia2024taskscaling,zhang2026prescriptive}. Our log-sigmoid environment learning curves in Section~\ref{sec:emerging-scaling} are closest in mathematical form to this bounded-performance line of work, but the independent variable is elapsed environment interaction within a task rather than pretraining compute or model scale.

Test-time and inference-time scaling provide a second point of comparison: different test-time methods exhibit their own empirical scaling behavior. One line scales plain repeated sampling: Evaluating Large Language Models Trained on Code introduced pass@\(k\) as a repeated-sampling evaluation for code generation~\cite{chen2021humaneval}, AlphaCode used large-scale sampling, filtering, and selection to improve competitive-programming solve rates~\cite{li2022alphacode}, and Large Language Monkeys showed that repeated sampling can scale coverage across orders of magnitude when outputs are automatically verifiable~\cite{brown2024largelanguagemonkeys}. A second line studies how test-time compute should be allocated across search, revision, voting, and verifier- or reward-guided strategies~\cite{snell2024testtime,wu2024inference}. A third line scales long-chain-of-thought inference in reasoning models: OpenAI's o1 report made train-time reinforcement learning compute and test-time thinking compute explicit scaling axes~\cite{openai2024learningreason}, and DeepSeek-R1 studies how reinforcement learning can elicit reasoning behaviors in LLMs~\cite{deepseekai2025deepseekr1}. These studies mostly scale the compute spent producing or selecting answers on traditional math, code, proof, or game-style tasks, and the reported functional forms are often local log-linear trends, exponentiated power laws, or frontiers for allocating compute. \benchmark{} instead scales \emph{interaction time}: it measures how an agent's best-so-far performance changes as elapsed time in an executable environment increases and as the agent repeatedly reads, acts, receives feedback, and revises. This gives broader task coverage and a different scaling object than sampling from a static prompt or selecting among answers.

Recent agentic test-time scaling work moves closer to our setting by allowing agents to acquire information from an environment over time. OpenAI's CUA report observed that computer-use performance improves when more steps are allowed~\cite{openai2025cua}, and BrowseComp reports smoother gains as browsing effort increases~\cite{openai2025browsecomp}. Thinking vs.\ Doing explicitly frames interaction length as a test-time scaling axis for web agents~\cite{shen2025thinkingdoing}. Closest in spirit, Mang et al.~\cite{mang2026humansstillbeatai} compare agents and humans on a long-horizon coding contest setting where agents can try, observe, and revise over time. They report an informative human reference curve and find that current agents plateau while strong humans continue improving. However, the measurement is concentrated in one contest-style domain and a limited task set. The human improvement curve is also difficult to extrapolate: the observed segment appears closer to linear improvement over time, but it may cover only the early portion of a longer sigmoid-like trajectory.

Scaling has also been studied in reinforcement learning. Hilton et al.~\cite{hilton2023rlscaling} introduced intrinsic performance to obtain smooth power-law relations between environment interactions, model size, and RL performance. Recent LLM RL scaling work fits sigmoidal compute-performance curves for post-training and uses smaller runs to predict larger RL runs~\cite{khatri2025scalingrlcompute}. Both RL and our evaluation measure learning from environment feedback rather than from human-labeled data. The practical difference is resolution and coverage. RL learning runs are expensive, so scaling evidence is usually gathered on narrower task distributions, fewer distinct environments, and fewer long trajectories. In \benchmark{}, the learning algorithm is in-context learning (ICL): the agent absorbs observations, diagnostics, and prior attempts into its context and uses them to improve subsequent actions. Because ICL is cheaper to repeat than RL training, \benchmark{} can measure environment learning across 134 executable task environments, multiple domains, full time-aligned trajectories, and repeated trials, yielding substantially broader environment coverage than is typical in RL scaling studies. The resulting fits go beyond a bounded-performance analogy to RL scaling: they provide a higher-resolution measurement of how agents learn from interaction.
